
\frontslide{Le hachage statique}

%------------------------------------------------------------
\section{Le hachage}
%------------------------------------------------------------
\begin{frame}{Le hachage}

Très utilisé comme structure en mémoire RAM. 

\vfill

Pour les bases de données, un concurrent de l'arbre-B
\begin{redbullet}
  \item \alert{Meilleur} (un peu, et en théorie) pour les recherches par clé
  \item \alert{N'occupe aucune place}
\end{redbullet}

\vfill

Mais

\vfill

\begin{redbullet}
  \item Se réorganise difficilement
  \item Ne supporte pas les recherches par intervalle
\end{redbullet}
  

\end{frame}
%------------------------------------------------------------
\begin{frame}{Principe du hachage}

Le stockage est organisé en $N$ \alert{fragments} (\textit{buckets})
constitués de séquences de blocs.

\vfill

La répartition des enregistrements se fait par un \alert{calcul}. Une \textit{fonction de hachage} $h$ prend une
         valeur de clé en entrée et retourne une adresse de fragment en sortie



\vfill
\begin{redbullet}
  \item \alert{Stockage}\,: on calcule l'adresse du fragment d'après la clé, on y stocke l'enregistrement
          
            \item \alert{Recherche (par clé)}\,:  on calcule l'adresse du fragment d'après la clé,
                on y cherche l'enregistrement
     
\end{redbullet}

\vfill

Simple et efficace\,!

\end{frame}
%------------------------------------------------------------
\begin{frame}{Exemple}

On veut créer une structure de hachage pour nos 16 films

\vfill

Simplifions\,: chaque fragment fait un bloc et contient 4 enregistrements au plus


\vfill

\begin{redbullet}
  \item on alloue 5 fragments (pour garder une marge de man{\oe}uvre)
  \item un \alert{répertoire} à 5 entrées (0 à 4) pointe vers les fragments
   \item On définit la fonction $h(titre) = rang(titre[0])\, mod\, 5$
\end{redbullet}

\vfill
\begin{blocgauche}
Donc on prend la première lettre du titre (par exemple 'I' pour 'Impitoyable'), on
prend sot rang dans l'alphabet (ici 9) et on garde le reste de la division par 5, le nombre de fragments.
\smallskip

$$h('Impitoyable')= 4$$
\end{blocgauche}

\end{frame}
%------------------------------------------------------------
\begin{frame}{Le résultat}

  \figSlideWithSize{hachage}{12cm}

\vfill
\begin{blocgauche}
La seule structure additionnelle est le répertoire, qui \alert{doit} tenir en RAM
\end{blocgauche}
\end{frame}
%------------------------------------------------------------
\begin{frame}[fragile]{Recherches}

\begin{redbullet}
  \item Par clé\,, Oui
  
 
\begin{minted}[baselinestretch=.9,
               fontsize=\small,
               framesep=2mm]{sql}
SELECT * FROM  Film WHERE titre = 'Impitoyable'
\end{minted}
 
   \item Par préfixe\,? Non

\begin{minted}[baselinestretch=.9,
               fontsize=\small,
               framesep=2mm]{sql}
SELECT * FROM  Film WHERE titre LIKE 'Mat%'
\end{minted}
 
   \item Par intervalle\,: non\,!

\begin{minted}[baselinestretch=.9,
               fontsize=\small,
               framesep=2mm]{sql}
SELECT *
FROM  Film
WHERE titre BETWEEN 'Annie Hall' AND 'Easy Rider'
\end{minted}
 
\end{redbullet}

\end{frame}
%------------------------------------------------------------
\begin{frame}{Les difficultés}


Très important\,: $h$ doit répartir uniformément les enregistrements
            dans les $n$ fragments
            
\vfill

Notre fonction est un contre-exemple\,: si une majorité de films commence par une
même lettre ('L' par exemple) la répartition va être déséquilibrée.

\vfill

\alert{Plus grave}\,: La structure simple décrite précédemment n'est pas
\alert{dynamique}.

\begin{redbullet}
  \item On ne peut pas changer un enregistrement de place

   \item Donc  il faut créer un chaînage de blocs quand
             un fragment déborde

  \item Et donc les performances se dégradent...
\end{redbullet}

\vfill

\begin{blocgauche}
Notez que l'on ne peut pas changer le nombre de fragments
car cela implique le changement de la fonction de répartition. 
\end{blocgauche}
\end{frame}
%------------------------------------------------------------
\begin{frame}{Exemple\,: insertion de Citizen Kane}

  \figSlideWithSize{hachage2}{12cm}

\vfill

\begin{blocgauche}
Il faut maintenant deux lectures pour accéder à Citizen Cane.
\end{blocgauche}
\end{frame}


%------------------------------------------------------------
\begin{frame}{Résumé: avantages et inconvénients du hachage}


En terme d'efficacité, il paraît optimal

\begin{redbullet}
  \item La structure (le répertoire) prend très peu de place
  \item Le coût d'une recherche par clé est constant
  \item De plus il est particulièrement simple à implanter.
\end{redbullet}

\vfill

Mais

\vfill
\begin{redbullet}
  \item Dans sa version de base, il se dégrade quand la situation évolue
  \item Pas de recherche pas intervalle ou préfixe
  \item Structure plaçante\,: il ne peut y avoir qu'une seule structure de hachage par table.
\end{redbullet}

\end{frame}
